(1) Snell 定律给 $n_1\sin\theta=n_2\sin\theta_t$。只有光从高折射率向低折射率入射且 $\sin\theta>n_2/n_1$，透射角才不能保持实数传播方向。临界角 $\theta_c=\arcsin(n_2/n_1)$，对 $\theta<\theta_c$ 仍有真实折射波，达到临界时其法向传播波矢趋零。
(2) 令真空波数 $k_0=2\pi/\lambda_0$，界面平行分量固定为 $k_x=k_0n_1\sin\theta$。介质2色散关系给 $k_{z2}^2=k_0^2n_2^2-k_x^2<0$，取 $k_{z2}=i\kappa$，$\kappa=k_0\sqrt{n_1^2\sin^2\theta-n_2^2}>0$。透射电场沿离界面法向 $z>0$ 按 $e^{-\kappa z}$ 衰减，幅度穿透深度 $d=1/\kappa$；强度的 $1/e$ 深度是 $1/(2\kappa)$。
(3) TE偏振下法向波矢 $k_{z1}=k_0n_1\cos\theta$ 为正实数，切向电场与磁场边界条件给 $r_{\rm TE}=(k_{z1}-i\kappa)/(k_{z1}+i\kappa)$。分子、分母互为共轭，故 $|r_{\rm TE}|=1$；相角为 $\phi_{\rm TE}=-2\arctan(\kappa/k_{z1})$，与几何反射相位不同。若做多次全反射，这个角度相关相位会影响偏振态。
(4) 单个无损倏逝平面波的时间平均法向坡印廷分量为零，因电磁场的有关正交分量存在 $90^\circ$ 相移；但是场强和沿界面的能流可非零，因此不能称介质2完全没有电磁能。若在衰减深度内放置第三个光学介质，两个界面的倏逝场可耦合成传输波，即受抑全反射，此时第一界面反射率不再必然为一。
